%!TEX root = main.tex
\chapter*{Abstract}
\addcontentsline{toc}{chapter}{Abstract}

Most published driver-distraction detection systems are trained and tested on images from one fixed camera position, in daylight, with speed usually measured only on a laptop, not real hardware. This project asks what happens when each of those three assumptions is removed, and answers each one with real measurements instead of estimates.

Three CNN architectures, a custom CNN, MobileNetV2 and EfficientNet-B0, were trained on the State Farm Distracted Driver Detection dataset. The training used a subject-wise split that holds entire drivers out of training, which avoids the subject leakage caused by a random image-level split. MobileNetV2 was chosen as the best model, with 84.67\% test accuracy and 84.01\% macro F1. Two robustness tests were then run on real data. For camera view, State Farm was combined with the SAM-DD multi-view dataset, and a frozen-backbone model reached only 75.94\% mixed-view accuracy, uneven across views, and fine-tuning raised this to 91.76\%. For lighting, the day-trained model was tested on real night footage recorded under active infrared light, replacing an earlier plan to make synthetic infrared-like images, since a synthetic transform can only approximate real infrared data, not replace it. The model's accuracy collapsed to 19.53\%, with zero recall on eight of the ten classes. Training on real night data brought night accuracy back up to 71.13\%, at a cost of 3.5 points of daytime accuracy.

The fine-tuned model was quantised to INT8 TensorFlow Lite, reducing its size 8.00$\times$, from 20.78~MB to 2.60~MB, for a cost of 2.7 accuracy points. It was then deployed on a Xilinx ZC702 board, after fixing a CPU instruction-set mismatch in the only published runtime for that platform. On-device accuracy was 77.3\%, with 100 out of 100 prediction agreement against the reference model. Inference took about 9.15 seconds per image, roughly 150$\times$ slower than the laptop figure and the clearest limitation of this deployment. A separate FPGA alert-decision stage was verified in simulation against real model predictions, including real misclassifications, and synthesised to 18 LUTs and 13 flip-flops, with 7.276~ns of positive timing slack.

The results show three things. First, camera-view and lighting robustness have to be trained for on purpose, they do not happen by themselves. Second, the cost of embedded deployment falls almost entirely on the classifier and not on the logic around it. Third, measuring a robustness problem before fixing it is what makes the fix easy to understand.
